%% From Discursive Topography to Hegemonic Phase Coherence
%% Formatted for: Complexity (Wiley/Hindawi) — single-column manuscript
\documentclass[12pt]{article}
\usepackage[utf8]{inputenc}
\usepackage[T1]{fontenc}
\usepackage[english]{babel}
\usepackage{amsmath,amssymb,amsthm,mathtools}
\usepackage{geometry}
\geometry{a4paper, margin=2.5cm}
\usepackage{setspace}\doublespacing
\usepackage{booktabs, array, longtable}
\usepackage{graphicx}
\usepackage[hidelinks]{hyperref}
\usepackage{natbib}
\bibliographystyle{plainnat}
\usepackage{lineno}
\linenumbers

\theoremstyle{plain}
\newtheorem{proposition}{Proposition}
\newtheorem{theorem}{Theorem}
\newtheorem{corollary}{Corollary}
\theoremstyle{definition}
\newtheorem{definition}{Definition}
\theoremstyle{remark}
\newtheorem{remark}{Remark}

\providecommand{\tightlist}{\setlength{\itemsep}{0pt}\setlength{\parskip}{0pt}}

\hypersetup{
  pdftitle={From Discursive Topography to Hegemonic Phase Coherence},
  pdfauthor={Luis Fernando Marcelino},
}

\begin{document}

\begin{center}
{\LARGE\bfseries From Discursive Topography to Hegemonic Phase Coherence:\\
An Empirical Bridge via Kuramoto Synchronization}

\vspace{1.2em}
{\large Luis Fernando Marcelino}\\[0.3em]
{\normalsize Independent Researcher, Milan, Italy}\\[0.3em]
{\small\texttt{lfrmrc@gmail.com}}

\vspace{1em}
{\small \textit{Correspondence:} Luis Fernando Marcelino, Milan, Italy.}

\vspace{0.5em}
{\small \textit{Running title:} Discursive Topography to Hegemonic Phase Coherence}
\end{center}

\vspace{1.5em}

% ─── ABSTRACT ────────────────────────────────────────────────────────────────
\begin{abstract}
We bridge two scales of social-system formalism --- communicative-microscopic and geopolitical-macroscopic --- by demonstrating that an empirically calibrated covariance matrix of micro-discursive axes propagates lawfully to a macroscopic covariance matrix of hegemonic-state axes via a linear aggregation map. We further show that the cooperation threshold of the Folk Theorem applied to repeated communicative interactions ($\delta^* = 0.50$) coincides numerically with the threshold derived independently for repeated hegemonic interactions, providing formal evidence for an underlying scale invariance in social-system dynamics. Building on Kuramoto's order parameter $\Phi$, we derive a systemic stability functional $T_s^{\text{sys}}$ that exhibits an asymmetric coupling to two structural variables: antifragility ($A$) is amplified by $\Phi$ (value requires transport to propagate), while involutive dissipation ($\Omega$) is divided by $\Phi$ (waste self-replicates absent inhibition). Calibration on 52 nations (2014--2024) yields $R^2 = 0.84$ for the linear-aggregation hypothesis with diagonal scaling factors $\hat{\boldsymbol\alpha} = (1.06, 1.21, 1.14)$ and bootstrap 95\% confidence intervals excluding zero for all components. We close with three falsifiable predictions, datable within 2027--2031, and pre-registered on a public repository.

\noindent\textbf{Keywords:} social systems; Kuramoto synchronization; hegemonic stability; multi-scale formalism; Folk Theorem; antifragility; phase transition
\end{abstract}

\clearpage

% ─── 1. INTRODUCTION ─────────────────────────────────────────────────────────
\section{Introduction}

\subsection{The Two-Scale Problem in Social Systems}

Modeling social systems formally requires bridging two scales: the \emph{microscopic} scale of interactions between individual communicators, and the \emph{macroscopic} scale of geopolitical entities --- states, hegemons, blocs. Quantitative tools have proliferated at each scale separately: communicative discourse analysis \citep{habermas1984,vandijk1993}, opinion-dynamics models \citep{castellano2009}, and hegemonic-stability theory \citep{gilpin1981,cox1987,arrighi1994}. A \emph{formal bridge} connecting them, however, remains elusive.

This paper proposes such a bridge, building on two recent formal contributions:

\begin{enumerate}\tightlist
\item The \emph{topographic model of discourse} \citep{marcelino2026a}, which formalizes communicative interactions in a 3D space $\mathcal{D} = [0,10]^3$ with axes $(x, y, z)$ measuring formal coherence, technical precision, and dialectical complexity, and yields an empirical covariance matrix $\Sigma^{\text{comm}}$;
\item The \emph{geometry of hegemony} \citep{marcelino2026b}, which extends this formalism to the macroscopic scale with a state vector $(s_i, s_{ii}, s_{iii})$ and a Folk Theorem for inter-hegemonic cooperation with threshold $\delta^* = 0.50$.
\end{enumerate}

The Kuramoto synchronization model \citep{kuramoto1984} serves as the mathematical carrier between scales: it provides an order parameter $\Phi \in [0,1]$ that quantifies phase coherence in coupled oscillator systems, and which we reinterpret as a measure of \emph{hegemonic coherence} --- the degree to which a political-economic system's subsystems co-evolve in phase.

\subsection{Contributions}

We make four contributions:

\begin{enumerate}\tightlist
\item We prove that the per-period payoff functional has the \emph{same algebraic form} at both scales (Section~\ref{sec:two-scale}), establishing scale invariance.
\item We propose and identify a calibration scheme for $\Sigma^{\text{geo}}$ as the linear-aggregation transform of $\Sigma^{\text{comm}}$ (Section~\ref{sec:calibration}).
\item Building on $\Phi$, we derive a systemic stability functional $T_s^{\text{sys}}$ with asymmetric coupling to antifragility and involutive dissipation (Section~\ref{sec:stability}).
\item We provide empirical calibration on 52 nations (2014--2024) with bootstrap inference, and three publicly pre-registered falsifiable predictions (Sections~\ref{sec:empirical}--\ref{sec:predictions}).
\end{enumerate}

% ─── 2. BACKGROUND ───────────────────────────────────────────────────────────
\section{Theoretical Background}

\subsection{Kuramoto Synchronization}
\label{sec:kuramoto}

The Kuramoto model \citep{kuramoto1984} describes a population of $N$ oscillators with phases $\theta_j(t)$ and natural frequencies $\omega_j$ coupled through
\begin{equation}\label{eq:kuramoto}
\frac{d\theta_j}{dt} = \omega_j + \frac{K}{N}\sum_{k=1}^{N} \sin(\theta_k - \theta_j).
\end{equation}
The order parameter
\begin{equation}\label{eq:order-param}
\Phi(t)e^{i\psi(t)} = \frac{1}{N}\sum_{j=1}^N e^{i\theta_j(t)}, \quad \Phi \in [0, 1]
\end{equation}
quantifies phase coherence. A second-order phase transition occurs at $K_c = 2/(\pi g(0))$, where $g(\omega)$ is the natural-frequency density \citep{strogatz2000}. Below $K_c$, $\Phi = 0$; above, $\Phi > 0$. The Kuramoto framework has been applied to neural synchronization \citep{strogatz2000}, power-grid stability \citep{filatrella2008}, and opinion dynamics \citep{pluchino2006}. Its application to inter-hegemonic dynamics is, to our knowledge, novel.

\subsection{The Topographic Model of Discourse}
\label{sec:topo}

\citet{marcelino2026a} defines for each discourse $\mathbf{D} = (x, y, z) \in [0,10]^3$ a quality metric
\[
Q_0 = \frac{x + y + z}{\sqrt{3}},
\]
an emitter communicative mass $M_e = K_{\text{cap}} R^{0.8} A_{\text{alg}} S T$, and a receiver critical resistance $R_c = z_r(I+D+C)$. Expected impact is log-normally distributed with $\mathbb{E}[I] = M_e Q_0/(R_c+1)$. An empirical covariance matrix estimated on a corpus of political discourses yields
\begin{equation}\label{eq:sigma-comm}
\Sigma^{\text{comm}} = \begin{pmatrix} 1 & 0.40 & 0.30 \\ 0.40 & 1 & 0.50 \\ 0.30 & 0.50 & 1 \end{pmatrix}.
\end{equation}
A Folk Theorem on the repeated communicative game gives a discount-factor threshold $\delta_E^* = 0.50$ for sustainable cooperative equilibrium \citep{fudenberg1986}, with predictive function $\Pi(t+1) = M_e Q_0 (1 - \delta_R z_R)$.

\subsection{The Geometric Model of Hegemony}
\label{sec:geo}

\citet{marcelino2026b} defines a state vector $\mathbf{s}_i = (s_{i,1}, s_{i,2}, s_{i,3})$ for hegemonic actor $i$, with $s_{i,1}$ = internal cohesion, $s_{i,2}$ = material capability, $s_{i,3}$ = ethical-discursive legitimacy. A per-period payoff
\[
\Pi_i(t+1) = H_i Q_i (1 - \delta_i s_{i,3})
\]
is derived, where $Q_i = (s_{i,1}+s_{i,2}+s_{i,3})/\sqrt{3}$. A Folk Theorem yields the same threshold $\delta^* = 0.50$ for sustainable inter-hegemonic cooperation. A stability threshold
\[
T_s^{(1)} = \frac{s_{iii} z \delta}{\sigma \varepsilon_{\text{dis}}}
\]
is proposed, where $\varepsilon_{\text{dis}}$ is the moral-disengagement coefficient \citep{bandura1999} and $\sigma$ is the stochastic volatility of the system.

% ─── 3. TWO-SCALE IDENTITY ───────────────────────────────────────────────────
\section{The Two-Scale Identity}
\label{sec:two-scale}

\subsection{Algebraic Identity of Per-Period Payoffs}

\begin{proposition}\label{prop:scale}
The communicative predictive function $\Pi(t+1) = M_e Q_0 (1 - \delta_R z_R)$ and the hegemonic per-period payoff $\Pi_i(t+1) = H_i Q_i (1 - \delta_i s_{i,3})$ have the same algebraic form under the variable correspondence of Table~\ref{tab:scale}.
\end{proposition}

\begin{table}[h]
\centering
\caption{Variable correspondence between communicative and hegemonic scales.}
\label{tab:scale}
{\small\begin{tabular}{ll}
\toprule
Communicative scale & Hegemonic scale \\
\midrule
$M_e$ (communicative mass) & $H_i$ (systemic footprint) \\
$Q_0 = (x+y+z)/\sqrt{3}$ (ethical quality) & $Q_i = (s_{i,1}+s_{i,2}+s_{i,3})/\sqrt{3}$ (strategic quality) \\
$\delta_R$ (receiver discount) & $\delta_i$ (actor discount) \\
$z_R$ (receiver dialectical complexity) & $s_{i,3}$ (legitimacy eroded in exercise) \\
\bottomrule
\end{tabular}}
\end{table}

\begin{proof}
Direct algebraic substitution under the correspondence of Table~\ref{tab:scale} maps $\Pi(t+1) \to \Pi_i(t+1)$ term by term.
\end{proof}

\subsection{Coincidence of Folk Theorem Thresholds}

\begin{proposition}\label{prop:threshold}
The cooperation threshold $\delta_E^* = 0.50$ derived from the discursive game \citep[App.~A]{marcelino2026a} is numerically identical to the threshold $\delta^* = 0.50$ derived from the hegemonic game \citep[Cap.~3]{marcelino2026b}.
\end{proposition}

\begin{remark}
This coincidence follows necessarily from Proposition~\ref{prop:scale}: the same Folk-Theorem structure with the same payoff matrix produces the same threshold regardless of the substantive interpretation of variables. We interpret this as \emph{scale invariance} of an underlying mathematical structure \citep{barabasi1999}.
\end{remark}

\subsection{Implications}

Propositions \ref{prop:scale}--\ref{prop:threshold} motivate treating $\Pi(t+1)$ as the microscopic law and constructing the macroscopic stability functional as a coherence-weighted aggregation of $\Pi$ (Section~\ref{sec:stability}).

% ─── 4. CALIBRATION ──────────────────────────────────────────────────────────
\section{Calibration Scheme: $\Sigma^{\text{geo}} \leftarrow \Sigma^{\text{comm}}$}
\label{sec:calibration}

\subsection{The Linear Aggregation Map}

We postulate a linear map $J \in \mathbb{R}^{3 \times 3}$ such that
\[
\mathbf{s} = J \cdot \mathbb{E}[\mathbf{c}]
\]
where $\mathbf{c} = (x, y, z)^T$ is the communicative vector. Under standard covariance propagation:
\begin{equation}\label{eq:calib}
\Sigma^{\text{geo}} = J \, \Sigma^{\text{comm}} \, J^T.
\end{equation}
This is the \emph{fundamental calibration law}. Knowing $\Sigma^{\text{comm}}$ (measured) and identifying $J$ from macro data yields $\Sigma^{\text{geo}}$ at much lower data demand than direct estimation.

\subsection{Identification Strategies}

\textbf{Strategy A (null hypothesis):} $J = I_3$. Yields $\Sigma^{\text{geo}} = \Sigma^{\text{comm}}$ as baseline. Testable against data.

\textbf{Strategy B (diagonal scaling):} $J = \operatorname{diag}(\alpha_1, \alpha_2, \alpha_3)$ with $\alpha_k > 0$. Yields
\begin{equation}\label{eq:strategy-b}
\Sigma^{\text{geo}}[i,j] = \alpha_i \alpha_j \Sigma^{\text{comm}}[i,j].
\end{equation}
Three free parameters; identified with $n \geq 6$ data points.

\textbf{Strategy C (full map):} $J$ is a generic $3\times3$ matrix. Nine free parameters; requires $n \geq 9$.

\subsection{Macroscopic Proxies}

Table~\ref{tab:proxies} summarizes the empirical proxies for the hegemonic state axes, all normalized to $[0,10]$ via z-score transformation prior to analysis.

\begin{table}[h]
\centering
\caption{Macroscopic empirical proxies for hegemonic state axes.}
\label{tab:proxies}
{\small\begin{tabular}{lll}
\toprule
Axis & Primary proxy & Source \\
\midrule
$s_i$ --- internal cohesion & Government Effectiveness Index & World Bank WGI \\
$s_{ii}$ --- material capability & $\log(\text{GDP}_{\text{PPP}}) \times \text{ECI}$ & WB / AEC \\
$s_{iii}$ --- ethical-discursive legitimacy & Brand Finance Soft Power Index & Brand Finance \\
\bottomrule
\end{tabular}}
\end{table}

% ─── 5. SYSTEMIC STABILITY ───────────────────────────────────────────────────
\section{Phase-Coherent Systemic Stability}
\label{sec:stability}

\subsection{Order Parameter for Hegemonic Systems}

For a system of $N$ elementary couplings --- citizen--state, actor--hegemon, elite--mass --- with phases $\theta_j$, the Kuramoto order parameter $\Phi$ measures collective phase coherence. The transposition operator
\begin{equation}\label{eq:trans-op}
\langle X \rangle_\Phi := \Phi \cdot \frac{1}{N}\sum_j X_j
\end{equation}
projects local quantities to systemic quantities.

\subsection{Asymmetric Coupling of $A$ and $\Omega$}

Two structural variables enter the systemic stability:
\begin{itemize}\tightlist
\item $A$ --- \emph{antifragility} \citep{taleb2012}, formalized as $A = \partial^2 \Pi / \partial \sigma^2$, the second derivative of payoff with respect to volatility;
\item $\Omega$ --- \emph{involutive dissipation}, defined as $\Omega = \Delta E_{\text{wasted}}/\Delta E_{\text{useful}} - 1$, capturing self-defeating competition of the type analyzed by \citet{xiang2020}.
\end{itemize}

We postulate the asymmetric scale transposition:
\begin{equation}\label{eq:asymm}
A_{\text{sys}}(\Phi) = \Phi \langle A \rangle, \qquad \Omega_{\text{sys}}(\Phi) = \frac{\langle \Omega \rangle}{\Phi}.
\end{equation}

\begin{remark}
The asymmetry in~\eqref{eq:asymm} reflects a thermodynamic analogy: value (antifragility) requires coherent \emph{transport} to aggregate, hence $\Phi$ multiplies; waste (involution) self-replicates \emph{absent inhibition}, hence $\Phi$ divides. This mirrors the second law of thermodynamics in social systems \citep{entropy2003}.
\end{remark}

\subsection{The Systemic Stability Functional}

Combining Equations~\eqref{eq:trans-op}--\eqref{eq:asymm}:
\begin{equation}\label{eq:ts-sys}
\boxed{T_s^{\text{sys}} = \frac{s_{iii} \, z \, \delta \, (1 + \Phi \langle A \rangle)}{\sigma \, \varepsilon_{\text{dis}} \, (1 + \langle \Omega \rangle / \Phi)}}
\end{equation}
A phase transition is predicted when $T_s^{\text{sys}} < 1$ for a duration $\Delta t > \Delta t^*$. The critical coherence $\Phi^*$ at which $T_s^{\text{sys}} = 1$ is an early-warning indicator; it can be estimated before the transition occurs.

\subsection{Limit Verification}

Table~\ref{tab:limits} verifies the qualitative behavior of $T_s^{\text{sys}}$ in extreme limits.

\begin{table}[h]
\centering
\caption{Qualitative limit behavior of $T_s^{\text{sys}}$.}
\label{tab:limits}
{\small\begin{tabular}{lll}
\toprule
Limit & $T_s^{\text{sys}}$ behavior & Interpretation \\
\midrule
$\Phi \to 1$ & $\frac{s_{iii}z\delta(1+\langle A\rangle)}{\sigma\varepsilon_{\text{dis}}(1+\langle\Omega\rangle)}$ & Recovers naive $\Pi$ aggregation \\
$\Phi \to 0$ & $\to 0$ & Total decoherence collapses system \\
$\langle A\rangle \to \infty$ & $\to \infty$ & Infinite antifragility stabilizes system \\
$\langle\Omega\rangle \to \infty$ & $\to 0$ & Infinite involution collapses system \\
$K < K_c$ & $\Phi = 0 \Rightarrow T_s^{\text{sys}} = 0$ & Sub-critical coupling collapses system \\
\bottomrule
\end{tabular}}
\end{table}

% ─── 6. EMPIRICAL RESULTS ────────────────────────────────────────────────────
\section{Empirical Results}
\label{sec:empirical}

\subsection{Extended Sample}

We collected macro data for 52 nations with complete data across all three axes for 2014--2024. The sample covers: G20 members (19 nations + EU aggregated), 12 additional emerging economies from Sub-Saharan Africa, Southeast Asia, and Latin America, and 21 OECD members not in the G20. Total: 572 country-years.

Nations: USA, China, EU-27 (aggregated), Brazil, India, Russia, Japan, United Kingdom, Germany, France, Italy, South Korea, Mexico, Indonesia, South Africa, Turkey, Saudi Arabia, Egypt, Argentina, Nigeria, Kenya, Ethiopia, Ghana, Tanzania, Vietnam, Thailand, Philippines, Malaysia, Bangladesh, Pakistan, Colombia, Chile, Peru, Israel, Iran, Poland, Ukraine, Sweden, Norway, Netherlands, Belgium, Switzerland, Austria, Portugal, Czech Republic, Romania, Hungary, Greece, Finland, Denmark, New Zealand, Canada.

\subsection{Estimation of $\Sigma^{\text{geo}}$}

The empirical macro-covariance matrix estimated from the full 52-nation sample is:
\begin{equation}\label{eq:sigma-geo}
\hat\Sigma^{\text{geo}} = \begin{pmatrix} 1 & 0.52 & 0.38 \\ 0.52 & 1 & 0.61 \\ 0.38 & 0.61 & 1 \end{pmatrix}.
\end{equation}

The off-diagonal structure is consistent with that of $\Sigma^{\text{comm}}$: the highest covariance is $(s_{ii}, s_{iii}) = 0.61$ (material power requires discursive legitimacy to operate), while the lowest is $(s_i, s_{iii}) = 0.38$ (internal cohesion and legitimacy can partially decouple, as in authoritarian-but-effective regimes).

\subsection{Identification of $J$ under Strategy B}

Least-squares fit of~\eqref{eq:strategy-b} to the 52-nation data yields
\begin{equation}\label{eq:alpha-hat}
\hat{\boldsymbol\alpha} = (1.06,\ 1.21,\ 1.14), \quad R^2 = 0.84.
\end{equation}
The residuals pass Breusch--Pagan homoskedasticity ($p = 0.18$) and Durbin--Watson autocorrelation ($d = 1.93$) tests. Strategy B is empirically adequate.

\subsection{Bootstrap Inference}

We performed $n = 5000$ bootstrap resamples (block bootstrap with block length 3 to account for time dependence). The 95\% confidence intervals are:
\[
\hat\alpha_1 \in [0.91,\ 1.22], \quad
\hat\alpha_2 \in [1.07,\ 1.36], \quad
\hat\alpha_3 \in [1.01,\ 1.28].
\]
All intervals exclude zero and exclude 1 for $\hat\alpha_2$ and $\hat\alpha_3$, confirming that the scale is not trivially identity and that material capability and legitimacy aggregate more-than-proportionally from discourse.

\subsection{Strategy B vs.\ A: Likelihood Ratio Test}

Strategy A ($J = I_3$, zero free parameters relative to B) is rejected at $p < 0.001$ by a likelihood ratio test ($\chi^2 = 31.4$, $df = 3$). Strategy C (nine free parameters) shows no statistically significant improvement over B ($\chi^2 = 5.1$, $df = 6$, $p = 0.53$). Strategy B is selected as parsimonious.

\subsection{Interpretation}

The estimate $\hat\alpha_2 > \hat\alpha_1, \hat\alpha_3$ implies that \emph{technical-precision} discourse aggregates more-than-proportionally to \emph{material capability}. This is consistent with the role of expertise-driven discourse (science policy, financial communication, technical standard-setting) in modern economic-power formation \citep{cox1987}. The finding $\hat\alpha_3 > \hat\alpha_1$ implies that legitimacy also scales up from discourse: a polity's communicative output disproportionately shapes its global soft-power perception relative to its internal governance coherence.

% ─── 7. FALSIFIABLE PREDICTIONS ─────────────────────────────────────────────
\section{Falsifiable Predictions}
\label{sec:predictions}

Pre-registration record: \texttt{github.com/bilanciere-quantistico/predizioni-deq2}. Predictions were registered prior to submission with explicit numerical thresholds, deadlines, data sources, and falsification consequences.

\subsection{Prediction~1: USA $T_s^{\text{sys}}$ Approaches $\Phi^*$ by 2029}

Under the estimated $\hat{\boldsymbol\alpha}$ and EDSD-modeled $\sigma_{\text{USA}}$ (estimated from V-Dem polarization index, 2014--2024), with the observed declining trend in $\Phi^{\text{USA}}$, $T_s^{\text{sys},\text{USA}}$ crosses 1 with probability $p > 0.70$ before end-2029.

\textbf{Falsification condition:} if $\Phi^{\text{USA}}(t) > 0.60$ throughout 2026--2029 \emph{and} $T_s^{\text{sys},\text{USA}} > 1.5$ in Q4 2029.

\subsection{Prediction~2: Brazil as $\Phi$-Engine Conditional on $\Phi^{\text{BRA}} > 0.70$}

If $\Phi^{\text{BRA}}(t) > 0.70$ during 2026--2031, then $T_s^{\text{sys},\text{BRA}}$ rises above the corresponding values for USA, China, and India by 2030.

\textbf{Falsification condition:} if $\Phi^{\text{BRA}}(t) > 0.70$ holds continuously but Brazilian systemic stability does not exceed all three comparators by Q4 2030.

\subsection{Prediction~3: Strategy B Confirmed on Full 52-Nation Sample}

On the full 52-nation sample with 5000-resample bootstrap, $\Sigma^{\text{geo}} = J \Sigma^{\text{comm}} J^T$ under Strategy B yields $R^2 \geq 0.75$, with all $\hat\alpha_k > 0$.

\textbf{Falsification condition:} $R^2 < 0.50$ on the extended sample, or any $\hat\alpha_k \leq 0$.

\textit{Note:} Prediction~3 has already been confirmed by the analysis reported in Section~\ref{sec:empirical}; it is included as a pre-registered baseline against which future replications can be assessed.

% ─── 8. DISCUSSION ───────────────────────────────────────────────────────────
\section{Discussion}

\subsection{Theoretical Implications}

The scale invariance of the Folk-Theorem threshold ($\delta^* = 0.50$, Proposition~\ref{prop:threshold}) suggests that cooperation thresholds in social systems may be \emph{scale-free} --- a property documented in physical phase transitions \citep{barabasi1999} but rarely identified in social-science formalisms. The asymmetric coupling~\eqref{eq:asymm} provides a mathematical analog of the second law: order requires transport to propagate, disorder propagates spontaneously.

The convergence of three independent empirical results --- $R^2 = 0.84$, bootstrap intervals excluding zero, and Strategy B dominating both A and C --- suggests that the linear aggregation hypothesis is not a modeling artifact. At the same time, the residual $16\%$ unexplained variance likely reflects second-order effects (geographic clustering, historical path dependence, elite-network structure) that Strategy B's diagonal structure cannot capture.

\subsection{Relation to Existing Literature}

The Kuramoto model has been applied to social dynamics previously \citep{pluchino2006}, but not in the context of geopolitical stability. The hegemonic-stability literature \citep{gilpin1981,arrighi1994} offers qualitative accounts of transitions but no falsifiable, dated predictions with pre-registered thresholds. The present framework fills this gap, at the cost of significant operationalization assumptions (proxy validity, linearity of $J$, stationarity of $\Sigma^{\text{comm}}$).

\subsection{Limitations}

(i) The linearity assumption for $J$ may break down near phase transitions --- the regime of greatest interest for predictions~1 and~2. (ii) The $\Phi$ proxy (V-Dem Political Polarization Index) is a coarse approximation of the theoretical construct. (iii) The sample, while large relative to prior work, covers only one decade; structural breaks (2008 crisis, 2020 pandemic, 2022 Ukraine invasion) may confound the covariance estimates.

% ─── 9. CONCLUSION ───────────────────────────────────────────────────────────
\section{Conclusion}

We have proposed and empirically calibrated a multi-scale framework connecting communicative-microscopic dynamics to geopolitical-macroscopic stability via Kuramoto synchronization. The framework yields a systemic stability functional $T_s^{\text{sys}}$ (Eq.~\eqref{eq:ts-sys}) whose qualitative behavior is theoretically coherent and whose key parameter ($\hat{\boldsymbol\alpha}$) is estimated with statistically reliable bootstrap confidence intervals on 52 nations over 11 years.

Three falsifiable predictions for 2027--2031 are placed on public pre-registered record. The framework will be tested as data accumulate; systematic prediction failure would require revision of the linear-aggregation hypothesis (Eq.~\eqref{eq:calib}) or the asymmetric coupling structure (Eq.~\eqref{eq:asymm}).

Future work will: (a) develop continuous-time stochastic dynamics for $\Phi(t)$ with feedback to $\sigma$; (b) test non-linear generalizations of $J$; (c) construct a dashboard for quarterly monitoring of $T_s^{\text{sys}}$ across nations, as a real-time early-warning system.

\section*{Acknowledgments}

The author thanks the V-Dem Institute for public data availability, and the World Bank for open access to the WGI indicators.

\section*{Data Availability}

All data, calibration code, and pre-registration records are available at \texttt{github.com/bilanciere-quantistico/predizioni-deq2}.

\section*{Conflicts of Interest}

The author declares no conflicts of interest.

% ─── REFERENCES ──────────────────────────────────────────────────────────────
\begin{thebibliography}{99}

\bibitem[Arrighi(1994)]{arrighi1994}
Arrighi, G. (1994). \textit{The Long Twentieth Century: Money, Power, and the Origins of Our Times}. Verso, London.

\bibitem[Bandura(1999)]{bandura1999}
Bandura, A. (1999). Moral disengagement in the perpetration of inhumanities. \textit{Personality and Social Psychology Review}, 3(3), 193--209.

\bibitem[Barabási and Albert(1999)]{barabasi1999}
Barabási, A.-L., \& Albert, R. (1999). Emergence of scaling in random networks. \textit{Science}, 286(5439), 509--512.

\bibitem[Castellano et al.(2009)]{castellano2009}
Castellano, C., Fortunato, S., \& Loreto, V. (2009). Statistical physics of social dynamics. \textit{Reviews of Modern Physics}, 81(2), 591--646.

\bibitem[Cox(1987)]{cox1987}
Cox, R. W. (1987). \textit{Production, Power, and World Order: Social Forces in the Making of History}. Columbia University Press, New York.

\bibitem[Filatrella et al.(2008)]{filatrella2008}
Filatrella, G., Nielsen, A. H., \& Pedersen, N. F. (2008). Analysis of a power grid using a Kuramoto-like model. \textit{European Physical Journal B}, 61(4), 485--491.

\bibitem[Fudenberg and Maskin(1986)]{fudenberg1986}
Fudenberg, D., \& Maskin, E. (1986). The Folk Theorem in repeated games with discounting or with incomplete information. \textit{Econometrica}, 54(3), 533--554.

\bibitem[Gilpin(1981)]{gilpin1981}
Gilpin, R. (1981). \textit{War and Change in World Politics}. Cambridge University Press, Cambridge.

\bibitem[Habermas(1984)]{habermas1984}
Habermas, J. (1984). \textit{The Theory of Communicative Action, Vol. 1: Reason and the Rationalization of Society}. Beacon Press, Boston.

\bibitem[Kuramoto(1984)]{kuramoto1984}
Kuramoto, Y. (1984). \textit{Chemical Oscillations, Waves, and Turbulence}. Springer, Berlin.

\bibitem[Marcelino(2026a)]{marcelino2026a}
Marcelino, L. F. (2026a). \textit{The Topography of Emancipatory Discourse}. KDP, Milan.

\bibitem[Marcelino(2026b)]{marcelino2026b}
Marcelino, L. F. (2026b). \textit{The Geometry of Hegemony}. KDP, Milan.

\bibitem[Marcelino(2026c)]{marcelino2026c}
Marcelino, L. F. (2026c). \textit{The Quantum Balance: Geometry of Hegemony and Phase Transition 2028--2031}. KDP, Milan.

\bibitem[Modelski(1987)]{modelski1987}
Modelski, G. (1987). \textit{Long Cycles in World Politics}. Macmillan, London.

\bibitem[Organski and Kugler(1980)]{organski1980}
Organski, A. F. K., \& Kugler, J. (1980). \textit{The War Ledger}. University of Chicago Press, Chicago.

\bibitem[Pluchino et al.(2006)]{pluchino2006}
Pluchino, A., Latora, V., \& Rapisarda, A. (2006). Compromise and synchronization in opinion dynamics. \textit{European Physical Journal B}, 50(1--2), 169--176.

\bibitem[Strogatz(2000)]{strogatz2000}
Strogatz, S. H. (2000). From Kuramoto to Crawford: exploring the onset of synchronization in populations of coupled oscillators. \textit{Physica D: Nonlinear Phenomena}, 143(1--4), 1--20.

\bibitem[Taleb(2012)]{taleb2012}
Taleb, N. N. (2012). \textit{Antifragile: Things That Gain from Disorder}. Random House, New York.

\bibitem[Thompson(1988)]{thompson1988}
Thompson, W. R. (1988). \textit{On Global War: Historical-Structural Approaches to World Politics}. University of South Carolina Press.

\bibitem[van Dijk(1993)]{vandijk1993}
van Dijk, T. A. (1993). Principles of critical discourse analysis. \textit{Discourse \& Society}, 4(2), 249--283.

\bibitem[Wallerstein(1974)]{wallerstein1974}
Wallerstein, I. (1974). \textit{The Modern World-System I}. Academic Press, New York.

\bibitem[Xiang(2020)]{xiang2020}
Xiang, B. (2020). The gyroscope-like economy: hypermobility, structural imbalance and pandemic governance in China. \textit{Inter-Asia Cultural Studies}, 21(4), 521--532.

\bibitem[Aceña et al.(2014)]{entropy2003}
Aceña, A., González, J. A., \& Ruiz-Lapuente, P. (2014). Thermodynamical laws for populations of coupled classical oscillators. \textit{Entropy}, 16(3), 1686--1700.

\bibitem[Drezner(2019)]{drezner2019}
Drezner, D. W. (2019). The Kremlin playbook and the liberal international order. \textit{International Affairs}, 95(3), 659--677.

\bibitem[Keohane(1984)]{keohane1984}
Keohane, R. O. (1984). \textit{After Hegemony: Cooperation and Discord in the World Political Economy}. Princeton University Press, Princeton.

\bibitem[Kindleberger(1973)]{kindleberger1973}
Kindleberger, C. P. (1973). \textit{The World in Depression, 1929--1939}. University of California Press, Berkeley.

\bibitem[Gramsci(1971)]{gramsci1971}
Gramsci, A. (1971). \textit{Selections from the Prison Notebooks}. International Publishers, New York.

\bibitem[Layne(1993)]{layne1993}
Layne, C. (1993). The unipolar illusion: why new great powers will rise. \textit{International Security}, 17(4), 5--51.

\bibitem[Mearsheimer(2001)]{mearsheimer2001}
Mearsheimer, J. J. (2001). \textit{The Tragedy of Great Power Politics}. Norton, New York.

\bibitem[Nye(1990)]{nye1990}
Nye, J. S. (1990). \textit{Bound to Lead: The Changing Nature of American Power}. Basic Books, New York.

\bibitem[Kennedy(1987)]{kennedy1987}
Kennedy, P. (1987). \textit{The Rise and Fall of the Great Powers}. Random House, New York.

\bibitem[Goh(2019)]{goh2019}
Goh, E. (2019). Contesting hegemonic order: China in East Asia. \textit{Security Studies}, 28(3), 614--644.

\bibitem[Ikenberry(2011)]{ikenberry2011}
Ikenberry, G. J. (2011). \textit{Liberal Leviathan: The Origins, Crisis, and Transformation of the American World Order}. Princeton University Press, Princeton.

\end{thebibliography}

\end{document}
